\documentclass[11pt,a4paper,oneside]{article}

% -------------------------------------------------
% Basic packages
% -------------------------------------------------
\usepackage[T1]{fontenc}
\usepackage[utf8]{inputenc}
\usepackage[spanish,es-noquoting,es-nodecimaldot]{babel}
\usepackage{lmodern}
\usepackage{microtype}
\usepackage[a4paper,margin=1.18in,headheight=15pt]{geometry}
\usepackage{amsmath,amssymb,amsfonts,amsthm,mathtools}
\usepackage{bm}
\usepackage{enumitem}
\usepackage{booktabs}
\usepackage{array}
\usepackage{xcolor}
\usepackage{hyperref}
\usepackage{setspace}
\usepackage{fancyhdr}
\usepackage{titlesec}
\usepackage[most]{tcolorbox}
\usepackage[capitalize,noabbrev]{cleveref}

\hypersetup{
  colorlinks=true,
  linkcolor=blue!50!black,
  citecolor=red!55!black,
  urlcolor=teal!60!black,
  pdftitle={Probabilidad y Estadística para Negocios -- Ejercicios estilo CFA},
  pdfauthor={César Galindo}
}

\theoremstyle{definition}
\newtheorem{definition}{Definición}[section]
\newtheorem{example}[definition]{Ejemplo}

\definecolor{ihesblue}{RGB}{25,62,114}
\definecolor{ihesgold}{RGB}{160,120,40}
\pagestyle{fancy}
\fancyhf{}
\fancyhead[L]{\textsc{Probabilidad y Estadística para Negocios}}
\fancyhead[R]{\textsc{César Galindo}}
\fancyfoot[C]{\thepage}
\renewcommand{\headrulewidth}{0.35pt}
\renewcommand{\footrulewidth}{0pt}
\titleformat{\section}{\Large\bfseries\color{ihesblue}}{\thesection.}{0.65em}{}
\titleformat{\subsection}{\large\bfseries\color{ihesblue}}{\thesubsection.}{0.65em}{}
\titlespacing*{\section}{0pt}{1.9em}{1.05em}
\titlespacing*{\subsection}{0pt}{1.45em}{0.95em}
\tcbset{colback=white,colframe=ihesblue!65!black,boxrule=0.5pt,arc=1.5pt}
\setlength{\parskip}{0.85em}
\setstretch{1.13}
\setlength{\parindent}{0pt}

\newtcolorbox{ideaBox}[1][]{enhanced,breakable,colback=blue!2,colframe=ihesblue,
  borderline west={1.5pt}{0pt}{ihesblue},title=Idea central,fonttitle=\bfseries,#1}
\newtcolorbox{summaryBox}[1][]{enhanced,breakable,colback=gray!4,colframe=black!55,
  borderline west={1.5pt}{0pt}{black!70},title=Fórmulas clave,fonttitle=\bfseries,#1}
\newtcolorbox{warningBox}[1][]{enhanced,breakable,colback=red!2,colframe=red!55!black,
  borderline west={1.5pt}{0pt}{red!75!black},title=Error frecuente,fonttitle=\bfseries,#1}
\newtcolorbox{cfaBox}[1][]{enhanced,breakable,
  colback=violet!3!white,colframe=violet!50!black,
  borderline west={2pt}{0pt}{violet!65!black},
  title=Ejercicio estilo CFA,fonttitle=\bfseries\color{violet!80!black},
  attach boxed title to top left={yshift=-2mm,xshift=4mm},
  boxed title style={colback=violet!12,colframe=violet!50!black,boxrule=0.5pt},#1}
\newtcolorbox{solBox}[1][]{enhanced,breakable,
  colback=green!2!white,colframe=green!45!black,
  borderline west={2pt}{0pt}{green!55!black},
  title=Solución,fonttitle=\bfseries\color{green!40!black},
  attach boxed title to top left={yshift=-2mm,xshift=4mm},
  boxed title style={colback=green!10,colframe=green!45!black,boxrule=0.5pt},#1}

\title{\vspace{-1.5cm}\textbf{\Large Probabilidad y Estadística para Negocios}\\[2pt]\large Ejercicios estilo CFA}
\author{César Galindo\\ \small Universidad Panamericana \textperiodcentered\ Facultad de Empresariales, Ciudad UP}
\date{\small}

\begin{document}
\maketitle

\begin{ideaBox}
Veinte problemas de opción múltiple, en el formato del examen CFA Nivel~I
(\emph{Quantitative Methods}: probabilidad, valor esperado y varianza,
covarianza y correlación, riesgo y rendimiento de portafolios,
distribuciones binomial, uniforme y normal, muestreo, intervalos de
confianza, prueba de hipótesis, y sesgo/curtosis), aplicados a contextos de
inversión y finanzas corporativas. Cada uno incluye la solución completa,
con todos los números verificados numéricamente. Se incluye antes un
recordatorio breve de las fórmulas usadas.
\end{ideaBox}

\begin{summaryBox}
\textbf{Reglas de probabilidad.} Para eventos $A,B$:
\[
P(A\cup B) = P(A)+P(B)-P(A\cap B), \qquad P(B\mid A) = \frac{P(A\cap B)}{P(A)}.
\]
$A$ y $B$ son independientes si $P(B\mid A)=P(B)$, equivalentemente
$P(A\cap B)=P(A)P(B)$.

\textbf{Teorema de Bayes.} Si $S_1,\dots,S_n$ son escenarios que cubren
todos los casos posibles (mutuamente excluyentes y colectivamente
exhaustivos), y $E$ es un evento observado:
\[
P(S_i\mid E) = \frac{P(E\mid S_i)\,P(S_i)}{P(E)}, \qquad\text{con}\qquad
P(E)=\sum_{j=1}^nP(E\mid S_j)\,P(S_j)\quad\text{(regla de la probabilidad total).}
\]

\textbf{Valor esperado, varianza y covarianza.} Para una variable aleatoria
discreta $R$ con escenarios $r_i$ y probabilidades $p_i$:
\[
E(R)=\sum_ip_ir_i, \qquad \mathrm{Var}(R)=\sum_ip_i\big(r_i-E(R)\big)^2, \qquad \sigma_R=\sqrt{\mathrm{Var}(R)}.
\]
Para dos variables $R_A,R_B$ con los mismos escenarios:
\[
\mathrm{Cov}(R_A,R_B)=\sum_ip_i\big(r_{A,i}-E(R_A)\big)\big(r_{B,i}-E(R_B)\big).
\]

\textbf{Binomial.} Si $X\sim\mathrm{Bin}(n,p)$, $P(X=k)=\binom{n}{k}p^k(1-p)^{n-k}$.

\textbf{Normal estándar.} Si $R\sim N(\mu,\sigma^2)$, entonces
$Z=(R-\mu)/\sigma\sim N(0,1)$, y $P(R\le r)=\Phi\!\left(\frac{r-\mu}{\sigma}\right)$.

\textbf{Error estándar de la media.} Para una muestra de tamaño $n$ de una
población con desviación estándar $\sigma$:
$\mathrm{SE}(\bar X)=\sigma/\sqrt{n}$.

\textbf{Portafolio de dos activos.} Con pesos $w_A+w_B=1$ y correlación
$\rho_{A,B}$:
\[
E(R_p)=w_AE(R_A)+w_BE(R_B), \qquad
\sigma_p^2=w_A^2\sigma_A^2+w_B^2\sigma_B^2+2w_Aw_B\rho_{A,B}\sigma_A\sigma_B,
\]
y el coeficiente de correlación es $\rho_{A,B}=\mathrm{Cov}(R_A,R_B)/(\sigma_A\sigma_B)$.

\textbf{Razón de Sharpe.} $\mathrm{Sharpe}=(E(R)-R_f)/\sigma$: rendimiento
en exceso de la tasa libre de riesgo, por unidad de riesgo total.

\textbf{Intervalo de confianza para la media.} Con error estándar
$\mathrm{SE}(\bar X)$ y valor crítico $z_{\alpha/2}$ (o $t_{\alpha/2}$ con
muestra pequeña y $\sigma$ desconocida):
$\bar X\pm z_{\alpha/2}\,\mathrm{SE}(\bar X)$.

\textbf{Prueba de hipótesis sobre una media.} Estadístico de prueba
$=\dfrac{\bar X-\mu_0}{\mathrm{SE}(\bar X)}$; se rechaza $H_0$ si el
estadístico excede, en valor absoluto, el valor crítico.
\end{summaryBox}

\begin{warningBox}
$P(B\mid A)\ne P(A\mid B)$ en general: confundir estas dos probabilidades
condicionales es el error más común al aplicar el teorema de Bayes (y
aparece constantemente en preguntas de examen sobre pronósticos de
analistas: la probabilidad de un ``profit warning'' dado que hay recesión
\emph{no es} la probabilidad de recesión dado que hay un profit warning).
\end{warningBox}

\section{Probabilidad condicional e independencia}

\begin{cfaBox}
Un analista estima que la probabilidad de que la economía crezca el próximo
trimestre es $P(A)=0.40$, que una empresa determinada supere las
expectativas de utilidades es $P(B)=0.30$, y que ambos eventos ocurran
simultáneamente es $P(A\cap B)=0.15$. ¿Cuál es $P(B\mid A)$, y son $A$ y
$B$ independientes?
\begin{enumerate}[label=\Alph*.]
\item $P(B\mid A)=0.300$; sí son independientes.
\item $P(B\mid A)=0.375$; no son independientes.
\item $P(B\mid A)=0.150$; no son independientes.
\end{enumerate}
\end{cfaBox}
\begin{solBox}
\textbf{B.} $P(B\mid A)=P(A\cap B)/P(A)=0.15/0.40=0.375$. Como
$P(B\mid A)=0.375\ne P(B)=0.30$, los eventos \textbf{no} son
independientes: que la economía crezca hace más probable que la empresa
supere expectativas.
\end{solBox}

\section{Teorema de Bayes}

\begin{cfaBox}
Un analista asigna probabilidades a tres escenarios económicos para el
próximo año: expansión, $P(\text{Exp})=0.50$; estabilidad,
$P(\text{Est})=0.30$; recesión, $P(\text{Rec})=0.20$. La probabilidad de
que una empresa emita un ``profit warning'' (advertencia de utilidades
menores a lo esperado) bajo cada escenario es: $P(W\mid\text{Exp})=0.10$,
$P(W\mid\text{Est})=0.20$, $P(W\mid\text{Rec})=0.50$. La empresa emite un
profit warning. ¿Cuál es la probabilidad actualizada de que la economía
esté en recesión, $P(\text{Rec}\mid W)$?
\begin{enumerate}[label=\Alph*.]
\item 20.0\%
\item 23.8\%
\item 47.6\%
\end{enumerate}
\end{cfaBox}
\begin{solBox}
\textbf{C.} Primero, por la regla de la probabilidad total,
\[
P(W)=0.50(0.10)+0.30(0.20)+0.20(0.50)=0.05+0.06+0.10=0.21.
\]
Por el teorema de Bayes,
\[
P(\text{Rec}\mid W)=\frac{P(W\mid\text{Rec})\,P(\text{Rec})}{P(W)}=\frac{0.20\times0.50}{0.21}=\frac{0.10}{0.21}\approx0.476=47.6\%.
\]
La probabilidad \emph{prior} de recesión era 20\%; observar el profit
warning la actualiza hacia arriba, a casi 47.6\%, porque un warning es
mucho más probable bajo recesión (50\%) que bajo los otros escenarios.
\end{solBox}

\section{Valor esperado y varianza}

\begin{cfaBox}
Los rendimientos de una acción bajo tres escenarios económicos son: auge
(\emph{boom}), probabilidad 25\%, rendimiento 20\%; normal,
probabilidad 50\%, rendimiento 8\%; recesión, probabilidad 25\%,
rendimiento $-10$\%. ¿Cuáles son el rendimiento esperado $E(R)$ y la
desviación estándar $\sigma_R$?
\begin{enumerate}[label=\Alph*.]
\item $E(R)=6.0\%$; $\sigma_R=10.00\%$
\item $E(R)=6.5\%$; $\sigma_R=10.71\%$
\item $E(R)=6.5\%$; $\sigma_R=114.75\%$
\end{enumerate}
\end{cfaBox}
\begin{solBox}
\textbf{B.}
\[
E(R)=0.25(20)+0.50(8)+0.25(-10)=5+4-2.5=6.5\%.
\]
\[
\mathrm{Var}(R)=0.25(20-6.5)^2+0.50(8-6.5)^2+0.25(-10-6.5)^2=0.25(182.25)+0.50(2.25)+0.25(272.25)=114.75,
\]
en unidades de (puntos porcentuales)$^2$; así,
$\sigma_R=\sqrt{114.75}\approx10.71\%$. (La opción C confunde la
varianza, en unidades al cuadrado, con la desviación estándar --- un error
común: la varianza \emph{nunca} se reporta directamente en las mismas
unidades que $R$.)
\end{solBox}

\section{Covarianza y diversificación}

\begin{cfaBox}
Dos escenarios igualmente relevantes para un par de activos $A$ y $B$:
escenario 1 (probabilidad 60\%), $R_A=10\%$, $R_B=5\%$; escenario 2
(probabilidad 40\%), $R_A=-2\%$, $R_B=8\%$. ¿Cuál es
$\mathrm{Cov}(R_A,R_B)$, y qué implica su signo para la diversificación de
un portafolio que combine ambos activos?
\begin{enumerate}[label=\Alph*.]
\item $\mathrm{Cov}=+8.64$; los activos se mueven juntos, poca diversificación.
\item $\mathrm{Cov}=-8.64$; los activos se mueven en direcciones opuestas, buena diversificación.
\item $\mathrm{Cov}=-8.64$; la covarianza negativa implica que no se debe invertir en ambos.
\end{enumerate}
\end{cfaBox}
\begin{solBox}
\textbf{B.} $E(R_A)=0.6(10)+0.4(-2)=5.2\%$,
$E(R_B)=0.6(5)+0.4(8)=6.2\%$. Entonces
\[
\mathrm{Cov}(R_A,R_B)=0.6(10-5.2)(5-6.2)+0.4(-2-5.2)(8-6.2)=0.6(-5.76)+0.4(-12.96)=-3.456-5.184=-8.64.
\]
Una covarianza negativa significa que, cuando $R_A$ sube respecto a su
media, $R_B$ tiende a bajar respecto a la suya (y viceversa): combinar
ambos activos reduce el riesgo total del portafolio por debajo del
promedio ponderado de sus riesgos individuales --- exactamente el
mecanismo detrás de la diversificación. (Nótese que esto \textbf{no}
implica evitar el activo, como sugiere C: covarianza negativa es
\emph{deseable} en un portafolio.)
\end{solBox}

\section{Distribución binomial}

\begin{cfaBox}
Un gestor de fondos supera a su índice de referencia (\emph{benchmark})
con probabilidad 60\% en cualquier trimestre dado, de forma independiente
trimestre a trimestre. ¿Cuál es la probabilidad de que supere al
benchmark en exactamente 3 de los próximos 4 trimestres?
\begin{enumerate}[label=\Alph*.]
\item 25.92\%
\item 34.56\%
\item 60.00\%
\end{enumerate}
\end{cfaBox}
\begin{solBox}
\textbf{B.} Con $X\sim\mathrm{Bin}(4,0.6)$,
\[
P(X=3)=\binom{4}{3}(0.6)^3(0.4)^1=4\times0.216\times0.4=0.3456=34.56\%.
\]
\end{solBox}

\section{Distribución normal}

\begin{cfaBox}
El rendimiento anual de un fondo se modela como
$R\sim N(\mu=12\%,\sigma=8\%)$. ¿Cuál es, \textbf{aproximadamente}, la
probabilidad de que el rendimiento del próximo año sea negativo?
\begin{enumerate}[label=\Alph*.]
\item 6.68\%
\item 15.87\%
\item 50.00\%
\end{enumerate}
\end{cfaBox}
\begin{solBox}
\textbf{A.} Estandarizando, $Z=(0-12)/8=-1.5$. De la tabla normal
estándar, $\Phi(-1.5)\approx0.0668$, es decir $P(R<0)\approx6.68\%$.
(La opción B, $15.87\%$, es en realidad $P(R>20\%)$ para este mismo
fondo, ya que $Z=(20-12)/8=+1$ y $1-\Phi(1)\approx0.1587$: un distractor
típico de examen que intercambia la cola.)
\end{solBox}

\section{Razón de seguridad de Roy (\emph{safety-first ratio})}

\begin{cfaBox}
Un inversionista requiere que su portafolio no rinda menos de
$R_L=2\%$. El Portafolio A tiene $E(R_A)=10\%$, $\sigma_A=10\%$; el
Portafolio B tiene $E(R_B)=8\%$, $\sigma_B=6\%$. Según el criterio de
seguridad de Roy, ¿cuál portafolio minimiza la probabilidad de rendir por
debajo de $R_L$?
\begin{enumerate}[label=\Alph*.]
\item El Portafolio A, porque tiene mayor rendimiento esperado.
\item El Portafolio B, porque tiene la razón SF más alta.
\item Son equivalentes, porque ambos tienen la misma prima de riesgo relativa.
\end{enumerate}
\end{cfaBox}
\begin{solBox}
\textbf{B.} La razón de Roy es
$\mathrm{SF}=\dfrac{E(R)-R_L}{\sigma}$; \emph{mayor} SF implica \emph{menor}
probabilidad de caer por debajo de $R_L$ (bajo normalidad). Aquí,
\[
\mathrm{SF}_A=\frac{10-2}{10}=0.80, \qquad \mathrm{SF}_B=\frac{8-2}{6}=1.00.
\]
Aunque A tiene mayor rendimiento esperado, B tiene la razón SF más alta
(por su menor volatilidad) y por lo tanto la menor probabilidad de un
rendimiento inferior al umbral: el criterio de Roy explícitamente
\textbf{no} elige por rendimiento esperado solo, sino por esa probabilidad
de cola.
\end{solBox}

\section{Muestreo y error estándar}

\begin{cfaBox}
La desviación estándar poblacional de los rendimientos diarios de un
activo es $\sigma=1.5\%$. Un analista toma una muestra de $n=36$ días de
operación. ¿Cuál es el error estándar de la media muestral?
\begin{enumerate}[label=\Alph*.]
\item 0.042\%
\item 0.250\%
\item 1.500\%
\end{enumerate}
\end{cfaBox}
\begin{solBox}
\textbf{B.} $\mathrm{SE}(\bar X)=\sigma/\sqrt{n}=1.5/\sqrt{36}=1.5/6=0.25\%$.
Cuadruplicar el tamaño de muestra ($n\to4n$) reduciría el error estándar a
la mitad, no a un cuarto: el error estándar decrece con $\sqrt n$, no con
$n$ --- un punto que aparece con frecuencia en preguntas de examen sobre
el efecto de aumentar el tamaño de muestra.
\end{solBox}

\section{Reglas de probabilidad: unión de eventos}

\begin{cfaBox}
En una encuesta a 200 empresas del IPC, 80 reportaron haber adoptado
inteligencia artificial en sus procesos ($P(A)=0.40$), 60 reportaron haber
tenido una ronda de financiamiento en el último año ($P(B)=0.30$), y 24
reportaron ambas cosas ($P(A\cap B)=0.12$). ¿Cuál es la probabilidad de
que una empresa elegida al azar haya adoptado IA \textbf{o} haya tenido
financiamiento (o ambas)?
\begin{enumerate}[label=\Alph*.]
\item 58\%
\item 70\%
\item 82\%
\end{enumerate}
\end{cfaBox}
\begin{solBox}
\textbf{A.} Por la regla de adición,
\[
P(A\cup B)=P(A)+P(B)-P(A\cap B)=0.40+0.30-0.12=0.58=58\%.
\]
Restar $P(A\cap B)$ evita contar dos veces a las empresas que hicieron
ambas cosas (una confusión frecuente que lleva a sumar $P(A)+P(B)$ sin
más, dando el distractor incorrecto de $70\%$).
\end{solBox}

\section{Distribuciones continuas: por qué lognormal, no normal}

\begin{cfaBox}
Al modelar el \textbf{precio} futuro de una acción (a diferencia de su
\emph{rendimiento}), los analistas suelen preferir una distribución
lognormal sobre una distribución normal. ¿Cuál es la razón
\textbf{principal}?
\begin{enumerate}[label=\Alph*.]
\item La distribución lognormal tiene media y varianza iguales.
\item Una variable normal puede tomar valores negativos, pero un precio no puede ser negativo.
\item La distribución lognormal es simétrica, igual que la normal.
\end{enumerate}
\end{cfaBox}
\begin{solBox}
\textbf{B.} Un precio de acción está acotado por abajo en cero (no puede
ser negativo), mientras que una variable normal tiene soporte en toda la
recta real $(-\infty,\infty)$ y por lo tanto asignaría probabilidad
positiva a precios negativos --- algo sin sentido económico. Si el
\emph{rendimiento compuesto continuamente} $\ln(P_1/P_0)$ se modela como
normal, entonces el precio mismo $P_1=P_0\,e^{\ln(P_1/P_0)}$ resulta
lognormal, que es siempre positivo por construcción. (C es falsa: la
lognormal está sesgada a la derecha, no es simétrica.)
\end{solBox}

\section{Regla de multiplicación: eventos independientes}

\begin{cfaBox}
Una empresa tiene dos proyectos de inversión totalmente independientes
entre sí, cada uno con 70\% de probabilidad de tener éxito. ¿Cuál es la
probabilidad de que \textbf{ambos} proyectos tengan éxito?
\begin{enumerate}[label=\Alph*.]
\item 35.0\%
\item 49.0\%
\item 70.0\%
\end{enumerate}
\end{cfaBox}
\begin{solBox}
\textbf{B.} Para eventos independientes, $P(A\cap B)=P(A)\,P(B)$. Aquí,
$0.70\times0.70=0.49=49.0\%$. (A es el promedio de las probabilidades, no
su producto; C repite la probabilidad individual --- ninguno de los dos es
la regla correcta para ``ambos ocurren''.)
\end{solBox}

\section{Rendimiento esperado de un portafolio}

\begin{cfaBox}
Un portafolio está compuesto en 60\% por acciones, con rendimiento
esperado $E(R_{\text{acc}})=12\%$, y en 40\% por bonos, con
rendimiento esperado $E(R_{\text{bon}})=5\%$. ¿Cuál es el rendimiento
esperado del portafolio, $E(R_p)$?
\begin{enumerate}[label=\Alph*.]
\item 8.50\%
\item 9.20\%
\item 10.20\%
\end{enumerate}
\end{cfaBox}
\begin{solBox}
\textbf{B.} El rendimiento esperado de un portafolio es el promedio
ponderado de los rendimientos esperados de sus componentes:
\[
E(R_p)=0.60(12)+0.40(5)=7.2+2.0=9.2\%.
\]
(A es el promedio simple, sin ponderar por los pesos 60/40 --- un error
frecuente.)
\end{solBox}

\section{Riesgo de un portafolio de dos activos}

\begin{cfaBox}
Un portafolio combina el Activo A ($w_A=0.5$, $\sigma_A=20\%$) y el
Activo B ($w_B=0.5$, $\sigma_B=10\%$), con correlación
$\rho_{A,B}=0.3$. ¿Cuál es la desviación estándar del portafolio,
$\sigma_p$?
\begin{enumerate}[label=\Alph*.]
\item 12.45\%
\item 15.00\%
\item 22.36\%
\end{enumerate}
\end{cfaBox}
\begin{solBox}
\textbf{A.}
\[
\sigma_p^2=w_A^2\sigma_A^2+w_B^2\sigma_B^2+2w_Aw_B\rho_{A,B}\sigma_A\sigma_B
=(0.5)^2(20)^2+(0.5)^2(10)^2+2(0.5)(0.5)(0.3)(20)(10),
\]
\[
\sigma_p^2=100+25+30=155, \qquad \sigma_p=\sqrt{155}\approx12.45\%.
\]
(B, $15\%$, sería el promedio ponderado \emph{de las desviaciones
estándar} --- correcto solo si $\rho_{A,B}=1$, es decir si los activos
estuvieran perfectamente correlacionados, lo cual anularía el beneficio de
diversificar. La correlación imperfecta, $\rho=0.3<1$, es precisamente lo
que hace que $\sigma_p$ sea \emph{menor} que ese promedio ponderado.)
\end{solBox}

\section{Coeficiente de correlación}

\begin{cfaBox}
Dos activos tienen $\mathrm{Cov}(R_A,R_B)=45$, $\sigma_A=10\%$ y
$\sigma_B=6\%$. ¿Cuál es su coeficiente de correlación $\rho_{A,B}$?
\begin{enumerate}[label=\Alph*.]
\item 0.45
\item 0.75
\item 4.50
\end{enumerate}
\end{cfaBox}
\begin{solBox}
\textbf{B.} $\rho_{A,B}=\mathrm{Cov}(R_A,R_B)/(\sigma_A\sigma_B)=45/(10\times6)=45/60=0.75$.
(C, $4.50$, olvida dividir entre $\sigma_A\sigma_B$; recuerda que $\rho$,
a diferencia de la covarianza, siempre está entre $-1$ y $+1$ --- un valor
de $4.50$ ya debería verse como imposible.)
\end{solBox}

\section{Distribución uniforme continua}

\begin{cfaBox}
El tiempo (en días) que tarda en cerrarse una ronda de financiamiento para
una \emph{startup}, dado que se inició el proceso, se modela como
$X\sim U(0,20)$. ¿Cuál es la probabilidad de que el cierre tome entre 5 y
15 días?
\begin{enumerate}[label=\Alph*.]
\item 25.0\%
\item 50.0\%
\item 75.0\%
\end{enumerate}
\end{cfaBox}
\begin{solBox}
\textbf{B.} Para $X\sim U(a,b)$, $P(c<X<d)=(d-c)/(b-a)$. Aquí,
$(15-5)/(20-0)=10/20=0.50=50.0\%$. En una uniforme, la probabilidad
depende solo de la \emph{longitud} del intervalo, no de su ubicación
dentro de $[a,b]$.
\end{solBox}

\section{Razón de Sharpe: comparación de fondos}

\begin{cfaBox}
El Fondo A tiene $E(R_A)=14\%$, $\sigma_A=18\%$; el Fondo B tiene
$E(R_B)=10\%$, $\sigma_B=8\%$. La tasa libre de riesgo es
$R_f=3\%$. Según la razón de Sharpe, ¿qué fondo tuvo mejor desempeño
ajustado por riesgo?
\begin{enumerate}[label=\Alph*.]
\item El Fondo A, con Sharpe $\approx0.611$.
\item El Fondo B, con Sharpe $=0.875$.
\item Ambos son iguales, porque el Fondo A tiene mayor rendimiento absoluto.
\end{enumerate}
\end{cfaBox}
\begin{solBox}
\textbf{B.}
\[
\mathrm{Sharpe}_A=\frac{14-3}{18}=\frac{11}{18}\approx0.611, \qquad
\mathrm{Sharpe}_B=\frac{10-3}{8}=\frac{7}{8}=0.875.
\]
A pesar de que A tiene mayor rendimiento esperado en términos absolutos, B
ofrece más rendimiento en exceso de $R_f$ \emph{por unidad de riesgo total}
--- esa es, por definición, la pregunta que responde la razón de Sharpe, y
B la responde mejor.
\end{solBox}

\section{Intervalo de confianza para la media}

\begin{cfaBox}
Una muestra de $n=100$ meses de rendimientos de un fondo tiene media
$\bar X=8\%$ y desviación estándar $\sigma=15\%$. Usando
$z_{0.025}=1.96$, ¿cuál es el intervalo de confianza al 95\% para el
rendimiento medio verdadero?
\begin{enumerate}[label=\Alph*.]
\item $(6.50\%,\ 9.50\%)$
\item $(5.06\%,\ 10.94\%)$
\item $(-7.00\%,\ 23.00\%)$
\end{enumerate}
\end{cfaBox}
\begin{solBox}
\textbf{B.} Primero, $\mathrm{SE}(\bar X)=\sigma/\sqrt n=15/\sqrt{100}=1.5\%$.
Luego,
\[
\bar X\pm z_{0.025}\,\mathrm{SE}(\bar X)=8\pm1.96(1.5)=8\pm2.94=(5.06\%,\ 10.94\%).
\]
(C usa $\sigma=15\%$ directamente en vez del error estándar de la
\emph{media} --- el error clásico de olvidar dividir entre $\sqrt n$ antes
de construir el intervalo.)
\end{solBox}

\section{Prueba de hipótesis sobre el alfa de un fondo}

\begin{cfaBox}
Un analista quiere probar si el alfa (exceso de rendimiento no explicado
por el mercado) de un fondo es significativamente distinto de cero,
$H_0:\alpha=0$ contra $H_a:\alpha\ne0$. La muestra tiene $n=64$
observaciones mensuales, alfa muestral $\hat\alpha=2\%$, y desviación
estándar muestral $s=8\%$. Al nivel de 5\% (valor crítico
$\approx1.96$ para una prueba de dos colas), ¿cuál es la conclusión?
\begin{enumerate}[label=\Alph*.]
\item Estadístico $=0.25$; no se rechaza $H_0$: el alfa no es significativo.
\item Estadístico $=2.00$; se rechaza $H_0$: el alfa es significativo.
\item Estadístico $=16.0$; no se puede concluir nada sin más información.
\end{enumerate}
\end{cfaBox}
\begin{solBox}
\textbf{B.} El error estándar es $\mathrm{SE}=s/\sqrt n=8/\sqrt{64}=8/8=1\%$.
El estadístico de prueba es
\[
t=\frac{\hat\alpha-0}{\mathrm{SE}}=\frac{2-0}{1}=2.00.
\]
Como $|2.00|>1.96$, se rechaza $H_0$ al 5\%: hay evidencia estadística de
que el alfa del fondo es distinto de cero. (A invierte la fracción,
$\mathrm{SE}/\hat\alpha$ en vez de $\hat\alpha/\mathrm{SE}$; C olvida
dividir entre $\sqrt n$.)
\end{solBox}

\section{Sesgo (\emph{skewness})}

\begin{cfaBox}
La distribución de pérdidas de una cartera de crédito tiene
media $>$ mediana $>$ moda, y una cola larga hacia la derecha (hacia
pérdidas grandes y poco frecuentes). Esta distribución se describe
\textbf{mejor} como:
\begin{enumerate}[label=\Alph*.]
\item Sesgada negativamente (a la izquierda).
\item Simétrica (sesgo cero).
\item Sesgada positivamente (a la derecha).
\end{enumerate}
\end{cfaBox}
\begin{solBox}
\textbf{C.} Cuando media $>$ mediana $>$ moda, y la cola larga apunta hacia
valores grandes (a la derecha), el sesgo es \textbf{positivo}: la media es
``jalada'' hacia arriba por unos pocos valores extremos grandes. En
distribuciones de pérdidas (o de rendimientos con riesgo de cola), un sesgo
positivo hacia pérdidas grandes es exactamente el patrón que preocupa a un
gestor de riesgo. (El orden se invierte, moda $>$ mediana $>$ media, en el
caso sesgado negativamente.)
\end{solBox}

\section{Curtosis y riesgo de cola}

\begin{cfaBox}
Los rendimientos diarios de un activo tienen curtosis (no en exceso) igual
a 7, comparado con 3 para una distribución normal. Un analista que asuma
erróneamente que estos rendimientos son normales estaría, \textbf{muy
probablemente}:
\begin{enumerate}[label=\Alph*.]
\item Sobreestimando la probabilidad de rendimientos extremos.
\item Subestimando la probabilidad de rendimientos extremos.
\item Sin ningún sesgo, porque la curtosis no afecta las colas de la distribución.
\end{enumerate}
\end{cfaBox}
\begin{solBox}
\textbf{B.} Una curtosis de 7 (mayor que 3) indica una distribución
\textbf{leptocúrtica}: más apuntada en el centro y con colas \emph{más
pesadas} (\emph{fat tails}) que la normal --- es decir, eventos extremos
son más frecuentes de lo que predice un modelo normal. Un analista que use
la normal de todas formas subestimará sistemáticamente la probabilidad de
pérdidas (o ganancias) grandes: precisamente el motivo por el que la
gestión de riesgo financiero rara vez confía únicamente en la normalidad.
\end{solBox}

\begin{summaryBox}
\textbf{En una frase:} en preguntas de probabilidad estilo CFA, el error
más costoso casi nunca es aritmético --- es usar la fórmula correcta con
el condicionamiento invertido (Bayes), confundir varianza con desviación
estándar, olvidar dividir entre $\sqrt n$ al construir un error estándar,
o elegir por rendimiento esperado ignorando el riesgo (Roy, Sharpe).
Verificar \emph{qué} se está condicionando, y en qué unidades está cada
resultado, antes de calcular.
\end{summaryBox}

\end{document}
